\documentclass[12pt,a4paper]{article}
\input{../pri-end/T0_preamble_standalone_De}

% Dokument-spezifische Metadaten
\hypersetup{
	pdftitle={Spin beeinflussen, auslesen und erweitern --- im FFGFT-Rahmen},
	pdfauthor={Johann Pascher},
	pdfsubject={FFGFT: Spin-Manipulation, Spin-Auslese und Qubits mit mehr als zwei Zuständen}
}

% Kopf-/Fußzeile für dieses Dokument
\fancyhead[L]{\small\textcolor{t0gray}{FFGFT · Spin und Qubit-Zustände}}
\fancyhead[R]{\small\textcolor{t0gray}{Johann Pascher · 2026}}

\title{%
	{\LARGE\bfseries\color{t0blue}
	 Spin beeinflussen, auslesen und erweitern}\\[0.5em]
	{\large Spin-Manipulation, Spin-Auslese und}\\[0.3em]
	{\large Qubits mit mehr als zwei Zuständen im FFGFT-Rahmen}\\[1em]
	{\normalsize Johann Pascher \quad $\cdot$ \quad März 2026}\\[0.5em]
	{\small \url{https://github.com/jpascher/T0-Time-Mass-Duality}}
}
\date{März 2026}

\begin{document}
	\maketitle
	\tableofcontents
	\newpage
	
	% ============================================================
% Inhalt: Spin lesen und schreiben --- und darüber hinaus
% Spin-Manipulation, Spin-Auslese und
% Qubits mit mehr als zwei Zuständen im FFGFT-Rahmen
% Neufassung März 2026
% ============================================================

% Dirac-Notation (braket-Paket nicht im Preamble)
\providecommand{\ket}[1]{\left|#1\right\rangle}
\providecommand{\bra}[1]{\left\langle#1\right|}
\providecommand{\braket}[2]{\left\langle#1\middle|#2\right\rangle}

% ============================================================
\section{Spin beeinflussen --- mehr als nur Magnetfelder}
\label{sec:spin-manipulieren}
% ============================================================

In Dokument~173 erschien der Spin im Rahmen der Ising-Maschine und
des minimalen Quantenresonators kurz: zwei erlaubte Orientierungen
einer Torsionskonfiguration --- Spin-up und Spin-down.
Die Frage wie man diese Orientierung von außen dreht
und wie man sie ausliest, ist die Brücke zwischen dem abstrakten
Resonatorbild und der technischen Realisierung.

Das Magnetfeld ist der historisch erste und anschaulichste Weg ---
aber bei weitem nicht der einzige.

% ============================================================
\subsection{Methode 1: Magnetfeld --- Zeeman-Effekt}
\label{sec:zeeman}
% ============================================================

Ein statisches Magnetfeld $\mathbf{B}$ spaltet die beiden
Spin-Zustände energetisch auf:
\begin{equation}
  \Delta E_Z \;=\; g^* \mu_B B,
\end{equation}
wobei $g^*$ der effektive $g$-Faktor des Materials und
$\mu_B$ das Bohr-Magneton sind.
Für GaAs gilt $g^* \approx -0{,}44$ ---
deutlich kleiner als beim freien Elektron ($g = 2$),
weil die Bandstruktur des Kristalls die effektive
Spin-Bahn-Kopplung modifiziert.

Ein resonanter Hochfrequenzpuls bei der Larmor-Frequenz
$\nu_L = g^* \mu_B B / h$ dreht den Spin gezielt um.
Das ist das Prinzip der Elektronenspinresonanz (ESR)
und --- im Kern --- auch der Magnetresonanztomographie.

\begin{keypoint}[FFGFT-Sprache: Magnetfeld als globales Torsionsfeld]
Ein äußeres Magnetfeld legt eine bevorzugte Richtung
im Torsionsfeld fest. Es bricht die Symmetrie zwischen
Spin-up und Spin-down energetisch auf.
Der Hochfrequenzpuls bei $\nu_L$ ist eine periodisch
oszillierende Torsionsstörung --- resonant, wenn ihre
Frequenz genau der Energiedifferenz der beiden
Torsionskonfigurationen entspricht.
\end{keypoint}

% ============================================================
\subsection{Methode 2: Elektrisches Feld --- EDSR}
\label{sec:edsr}
% ============================================================

Elektrische Felder koppeln nicht direkt an den Spin ---
das Elektron hat keine elektrische Dipolmoment-Komponente
im Spinraum. Aber indirekt koppeln sie sehr wohl,
über die \textbf{Spin-Bahn-Kopplung}:

Das Elektron, das sich im elektrischen Feld des Kristallgitters
bewegt, sieht in seinem eigenen Ruhesystem ein effektives
Magnetfeld. Dieses interne Feld dreht den Spin.
Ein schnell variierendes äußeres elektrisches Feld
kann diesen Mechanismus resonant ansteuern ---
das nennt sich \textbf{EDSR} (Electric Dipole Spin Resonance).

\begin{important}[Vorteil gegenüber Magnetfeldmethoden]
Elektrische Felder lassen sich mit Gate-Elektroden
auf wenigen Nanometern lokal erzeugen und in unter
einer Nanosekunde schalten --- Magnetfelder nicht.
Für integrierte Qubit-Arrays ist EDSR daher die
bevorzugte Methode zur Einzel-Qubit-Rotation.
\end{important}

In FFGFT-Sprache: das elektrische Feld deformiert
die Resonatorgeometrie lokal --- und über die
Spin-Bahn-Kopplung transliert diese Geometrieänderung
direkt in eine Torsionsänderung des Spinfreiheitsgrades.
Die Kopplung ist nicht direkt, aber geometrisch vermittelt.

% ============================================================
\subsection{Methode 3: Zirkulär polarisiertes Licht}
\label{sec:optisch}
% ============================================================

Zirkulär polarisiertes Licht ($\sigma^+$ oder $\sigma^-$)
trägt Drehimpuls $\pm\hbar$ pro Photon.
Bei resonanter Anregung eines Quantenpunkts überträgt
ein $\sigma^+$-Photon diesen Drehimpuls direkt auf
das erzeugte Exziton --- und damit auf den Elektronenspin.

Das erlaubt \textbf{rein optische Spinpräparation}
ohne jedes äußere Feld: durch Wahl der Photonpolarisation
wählt man den Spinzustand.

Ausgelesenes wird über die Polarisation des
Fluoreszenzlichts: je nach Spinzustand emittiert der
Quantenpunkt bevorzugt $\sigma^+$ oder $\sigma^-$-Licht.
Polarisationsmessung $\equiv$ Spinmessung.

\begin{foundation}[FFGFT: Photon überträgt Torsionsdrehimpuls]
In FFGFT trägt das Photon seinen Drehimpuls als
topologische Eigenschaft der elektromagnetischen
Feldtorsion. Bei der Absorption rotiert diese
Feldtorsion die Elektronentorsion in die
komplementäre Konfiguration.
Spin-up und Spin-down sind zwei entgegengesetzte
Torsionswindungen --- das zirkulare Photon
schaltet zwischen ihnen um, weil es selbst
eine gerichtete Windung trägt.
\end{foundation}

% ============================================================
\subsection{Methode 4: Austauschwechselwirkung}
\label{sec:exchange}
% ============================================================

Wenn zwei Elektronen in benachbarten Quantenpunkten
räumlich überlappen, koppeln ihre Spins direkt
aneinander --- die \textbf{Austauschwechselwirkung}
(exchange interaction) mit Stärke $J$:
\begin{equation}
  H_J \;=\; J(t)\;\mathbf{S}_1 \cdot \mathbf{S}_2.
\end{equation}
Durch Pulsieren des Gate-Feldes (das den Überlapp und
damit $J(t)$ steuert) kann man gezielte
Zwei-Qubit-Operationen ausführen.

Das ist die Basis für \textbf{Spin-Qubit-Gatter}
in Halbleitersystemen (Loss-DiVincenzo-Modell, 1998).
Die Stärke: keine externe Mikrowelle nötig ---
die Kopplung ist rein elektrisch steuerbar.

% ============================================================
\subsection{Methode 5: Hyperfein-Wechselwirkung}
\label{sec:hyperfein}
% ============================================================

Das Elektron im Quantenpunkt wechselwirkt mit den
Kernspins der $\sim 10^4$--$10^6$ umliegenden Atome.
Diese Kopplung ist normalerweise ein Problem:
das fluktuierende Kernspinbad erzeugt ein
zufälliges effektives Magnetfeld und zerstört
die Spinpolarisation in Zeiten von $T_2^* \sim 10$--100\,ns.

Aber kontrolliert eingesetzt ist sie ein Werkzeug:
durch dynamische Entkopplung (schnelle Pulsfolgen
die die Kernspin-Fluktuationen mitteln)
kann $T_2$ auf Mikrosekunden verlängert werden.
In einigen Systemen (z.\,B. Phosphor in Silizium)
wird der Kernspin selbst als hochstabiles Qubit genutzt
($T_1 > 1$\,s bei tiefen Temperaturen).

% ============================================================
\section{Spin auslesen --- fünf Wege}
\label{sec:spin-auslesen}
% ============================================================

% ============================================================
\subsection{Spin-zu-Ladungs-Konversion}
\label{sec:spin-ladung}
% ============================================================

Der wichtigste Weg in der Praxis.
Man wählt das Potential so, dass nur Spin-down
(oder nur Spin-up) energetisch erlaubt ist,
aus dem Quantenpunkt heraus zu tunneln.
Ein einzelnes tunnelndes Elektron erzeugt einen
messbaren Strompuls im benachbarten
\textbf{Einzelelektronen-Transistor} (SET).

Kein Signal: Spin-up (geblockt).
Signal: Spin-down (getunnelt).
Empfindlichkeit: einzelne Elektronen, Fehlerrate $< 1\,\%$.

\begin{keyresult}[Spin-zu-Ladungs-Konversion: technischer Stand 2025]
In Si/SiGe- und GaAs-Systemen erreichen
Spin-zu-Ladungs-Konversions-Readout-Schemen
Genauigkeiten $> 99{,}5\,\%$ bei Auslesezeiten
unter $1\,\mu$s --- ausreichend für fehlerkorrigierte
Quantencomputer nach dem Surface-Code-Schwellenwert.
\end{keyresult}

% ============================================================
\subsection{Pauli-Blockade}
\label{sec:pauli}
% ============================================================

Eleganter: kein externes Messgerät nötig.
Das Pauli-Prinzip selbst ist der Detektor.

Zwei benachbarte Quantenpunkte, je ein Elektron.
Man versucht beide Elektronen in denselben Punkt zu laden.
Sind die Spins antiparallel: geht (beide passen ins Grundniveau).
Sind die Spins parallel: geht nicht (Pauli-Verbot) ---
der Transport ist \textbf{blockiert}.

Messbar als Leitwert-Unterschied: Blockade $\equiv$ gleiche Spins.

\begin{foundation}[FFGFT: Pauli-Blockade als geometrische Inkompatibilität]
Zwei gleiche Torsionskonfigurationen können nicht
denselben Resonator belegen --- nicht wegen eines
postulierten Verbots, sondern weil dieselbe Torsionswindung
in einem Hohlraum keine zweite stabile Konfiguration
derselben Art zulässt.
Das Pauli-Prinzip ist in FFGFT eine geometrische Aussage
über die Kompatibilität von Feldkonfigurationen ---
keine unabhängige Quantisierungsregel.
\end{foundation}

% ============================================================
\subsection{Faraday- und Kerr-Rotation --- zustandserhaltende Auslese}
\label{sec:optisch-auslesen}
% ============================================================

Optische Auslese nutzt dass Spin-up und Spin-down
an verschiedene optische Übergänge koppeln.
Zirkulär polarisiertes Licht wird je nach Spinzustand
unterschiedlich gestreut --- die Polarisation
des Fluoreszenzlichts verrät den Spinzustand.
Diese resonante Fluoreszenz ist invasiv: das absorbierte
Photon regt das System an und verändert seinen Zustand.

Davon grundlegend verschieden ist die
\textbf{Faraday- bzw. Kerr-Rotation}:
Ein linear polarisierter Laserstrahl läuft am Quantenpunkt
\emph{vorbei} --- er wird nicht absorbiert, sondern dreht
seine Polarisationsebene um einen winzigen Winkel
$\theta_F \propto S_z$ beim Passieren.
Der Spinzustand ist aus dem Drehwinkel ablesbar,
ohne dass ein einziges Photon absorbiert wird.

\begin{keyresult}[Faraday-Rotation: Quantenzustand auslesen ohne ihn zu stören]
Die Faraday-Rotation ist eine
\textbf{Quantum-Non-Demolition-Messung} (QND) ---
eine Auslesemethode die den gemessenen Zustand
selbst nicht verändert.

Physikalisch: Das Lichtfeld koppelt dispersiv an den Spin ---
die Wechselwirkung ist so gewählt dass sie \emph{kommutiert}
mit dem Spin-Operator $S_z$:
\begin{equation}
  [H_\text{Faraday},\; S_z] \;=\; 0.
\end{equation}
Eine Observable die mit dem Hamiltonoperator kommutiert,
wird durch die Messung nicht gestört.
Die Polarisationsdrehung $\theta_F$ ist proportional zu $S_z$
--- der Eigenwert wird abgelesen, der Eigenzustand bleibt erhalten.

Zeitauflösung: im Pikosekundenbereich.
Empfindlichkeit: einzelne Spins in Quantenpunkten nachweisbar.
Der Spinzustand kann \emph{wiederholt} gemessen werden ---
jede Messung liefert dasselbe Ergebnis.
\end{keyresult}

\begin{foundation}[FFGFT: QND-Messung als Beweis für stabile Torsionskonfiguration]
Die Tatsache dass die Faraday-Rotation den Spinzustand
nicht verändert, ist in FFGFT kein Sonderfall ---
es ist die direkte Folge der Stabilität der Torsionskonfiguration.

Spin-up ist eine stabile topologische Torsion.
Ein vorbeilaufendes Photon dreht seine Polarisationsebene
weil es das Torsionsfeld passiert --- aber es bringt
die Torsionskonfiguration nicht aus ihrem stabilen Zustand,
weil die Kopplung so schwach und symmetrisch ist,
dass kein Energieübertrag in die Torsionsmode stattfindet.

Das ist dieselbe Logik wie bei einem Resonator:
ein schwaches Signal das die Resonanzfrequenz nicht trifft,
passiert den Resonator ohne ihn anzuregen.
Das Photon der Faraday-Messung trifft absichtlich
\emph{nicht} die Resonanzfrequenz des Übergangs ---
es liest die topologische Invariante ab, ohne sie zu stören.

\textbf{Der entscheidende Schluss:}
Wenn eine Messung denselben Zustand immer wieder liefert
und dabei den Zustand nicht verändert,
dann hatte der Zustand diesen Wert \emph{vor der Messung}.
Das ist der empirische Beweis dass stabile
Torsionskonfigurationen real und definiert existieren ---
unabhängig davon ob wir sie messen oder nicht.
Die QND-Messung beweist den Determinismus.
\end{foundation}

% ============================================================
\section{Die Grenze von zwei Zuständen --- und wie man sie überwindet}
\label{sec:jenseits-spin}
% ============================================================

Ein Spin-Qubit hat zwei Zustände: $\ket{\uparrow}$ und $\ket{\downarrow}$.
Das ist das absolute Minimum eines Quantenresonators
(Dokument~173, Abschnitt über den minimalen Resonator).

Aber ein Quantenpunkt ist kein Zwei-Zustands-System von Natur aus ---
er hat beliebig viele Energieniveaus.
Der Spin ist nur das \textit{einfachste} Merkmal.

Es gibt mindestens fünf Strategien wie man aus einem
Quantenpunkt mehr als zwei unterscheidbare Zustände
für Quanteninformation gewinnt.

% ============================================================
\subsection{Strategie 1: Qutrit und Qudit --- mehr Energieniveaus}
\label{sec:qudit}
% ============================================================

Statt $\ket{0}$ und $\ket{1}$ nutzt man $\ket{0}$, $\ket{1}$, $\ket{2}$
--- ein \textbf{Qutrit} (drei Zustände) ---
oder allgemein $\ket{0}$ bis $\ket{d-1}$:
ein \textbf{Qudit} der Dimension $d$.

Im Quantenpunkt entspricht das den ersten $d$ Energieniveaus.
Die Grundformel:
\begin{equation}
  E_n \;=\; \frac{n^2 \pi^2 \hbar^2}{2m_\text{eff} R^2},
  \quad n = 1, 2, \ldots, d.
\end{equation}
Für $d = 3$ (Qutrit) braucht man drei unterscheidbare Niveaus
mit $E_3 - E_2 \gg k_BT$ --- sonst verwischt die thermische
Besetzung die dritte Stufe.

\begin{important}[Bedingung für stabiles Qudit der Dimension $d$]
Alle $d$ Niveauabstände müssen die thermische Energie übersteigen:
\[
  E_{n+1} - E_n \;=\; (2n+1)\,E_1 \;\gg\; k_BT
  \quad \text{für alle } n = 1,\ldots,d-1.
\]
Der kritischste Abstand ist $E_2 - E_1 = 3E_1$ (der kleinste).
Qudits hoher Dimension erfordern kleine Quantenpunkte
oder tiefe Temperaturen --- oder beides.
\end{important}

Die Informationsdichte eines Qudits:
$\log_2(d)$ Bits pro Objekt.
Ein Qutrit kodiert $\log_2(3) \approx 1{,}585$ Bits ---
58\,\% mehr als ein Qubit pro Objekt.
Ein Qu\textit{fünf}-it kodiert $\log_2(5) \approx 2{,}32$ Bits.

\begin{keyresult}[Qudit-Effizienz auf der $\xi$-Stufe $N \approx 8$]
Ein Quantenpunkt mit $R = 3$\,nm bei 4\,K
hat mindestens 5 thermisch stabile Niveaus
(alle Abstände $\gg k_BT(4\,\text{K}) = 0{,}34$\,meV).
Das entspricht einem Qu\textit{fünf}-it mit
$\log_2(5) \approx 2{,}32$ Bits Informationsinhalt ---
in einem Resonator der 22\,nm Durchmesser hat.
\end{keyresult}

% ============================================================
\subsection{Strategie 2: Orbital-Qubit --- Geometrie statt Energie}
\label{sec:orbital}
% ============================================================

Ein Quantenpunkt hat nicht nur verschiedene Energieniveaus ---
er hat auch verschiedene \textbf{Orbitale} (räumliche
Wellenformuster mit gleicher Energie, aber unterschiedlicher
Geometrie): s-, p-, d-artige Zustände wie im Atom.

In einem nicht-sphärischen Quantenpunkt sind die
p-Orbitale energetisch aufgespalten.
Man kann die beiden untersten p-Zustände ($p_x$ und $p_y$)
als Qubit nutzen --- ein \textbf{Orbital-Qubit} ---
\emph{ohne} auf den Spin zu schauen.

Der Vorteil: Orbital-Zustände koppeln stark an elektrische
Felder, schwach an magnetische Störfelder.
Das macht sie potenziell kohärenter in magnetisch
verrauschten Umgebungen.

% ============================================================
\subsection{Strategie 3: Exziton-Zahl-Qudit}
\label{sec:exziton-qudit}
% ============================================================

Aus Dokument~173: ein Quantenpunkt kann $0, 1, 2, 3, \ldots$
Exzitonen gleichzeitig tragen --- diskrete Speicherstufen,
scharf unterscheidbar.

Diese Exzitonen-Zahl selbst ist ein Qudit:
\[
  \ket{0\,\text{Ex}},\quad
  \ket{1\,\text{Ex}},\quad
  \ket{2\,\text{Ex}},\quad
  \ket{3\,\text{Ex}},\;\ldots
\]
Das Auslesen erfolgt über die Emissionswellenlänge:
Ein Biexziton ($\ket{2\,\text{Ex}}$) emittiert bei einer
anderen Frequenz als ein Exziton ($\ket{1\,\text{Ex}}$)
--- spektral klar unterscheidbar.

\begin{foundation}[FFGFT: Exziton-Zahl als Torsionspaar-Zahl]
Jedes Exziton ist ein gebundenes Paar entgegengesetzter
Torsionskonfigurationen (Elektron + Loch = Torsion + Anti-Torsion).
Die Exziton-Zahl zählt wie viele solcher Paare
gleichzeitig im Resonator stabil koexistieren können.
Die Obergrenze ist durch die Resonatorgeometrie gegeben ---
zu viele Torsionspaare in einem zu kleinen Raum
führen zu instabilen Überlagerungen (Auger-Rekombination).
\end{foundation}

% ============================================================
\subsection{Strategie 4: Spin-Bahn-Qudit --- Spin und Orbital kombiniert}
\label{sec:spin-orbital}
% ============================================================

Der reichhaltigste Zustandsraum entsteht wenn man
\textbf{Spin} und \textbf{Orbital} kombiniert.

Ein Elektron in einem Quantenpunkt hat:
\begin{itemize}
  \item 2 Spinzustände ($\uparrow$, $\downarrow$)
  \item mehrere Orbitalzustände (s, $p_x$, $p_y$, \ldots)
\end{itemize}
Der kombinierte Zustandsraum:
\begin{equation}
  \ket{n_\text{orbital},\, m_\text{spin}}
  \;=\; \ket{1,\uparrow},\;\ket{1,\downarrow},\;
  \ket{2,\uparrow},\;\ket{2,\downarrow},\;\ldots
\end{equation}
Für die untersten zwei Orbitale ergibt das 4 Zustände ---
ein \textbf{Ququart} (Qudit der Dimension 4,
entspricht 2 verschränkten Qubits in einem einzigen Objekt).

Die Spin-Bahn-Kopplung verbindet die beiden Freiheitsgrade
und erlaubt es, durch elektrische Felder gezielt
zwischen den kombinierten Zuständen zu navigieren.

\begin{center}
\resizebox{\textwidth}{!}{%
\begin{tabular}{clcc}
\toprule
\textbf{Strategie} & \textbf{Freiheitsgrad} & \textbf{Dimension $d$} & \textbf{Bits pro QD} \\
\midrule
Spin-Qubit     & Spin           & 2      & 1{,}00 \\
Orbital-Qubit  & Orbital        & 2--4   & 1{,}00--2{,}00 \\
Qudit          & Energieniveau  & 3--6   & 1{,}58--2{,}58 \\
Exziton-Qudit  & Exziton-Zahl   & 3--5   & 1{,}58--2{,}32 \\
Spin-Orbital   & Spin + Orbital & 4--8   & 2{,}00--3{,}00 \\
\bottomrule
\end{tabular}
}
\end{center}

% ============================================================
\subsection{Strategie 5: Kontinuierliche Variablen ---
  der quantenmechanische Oszillator}
\label{sec:cv}
% ============================================================

Eine radikale Alternative zum diskreten Qudit-Konzept:
statt diskreter Zustände nutzt man
\textbf{kontinuierliche Variablen} (CV).

Ein harmonischer Oszillator hat unendlich viele
Energieniveaus $\ket{n}$, $n = 0, 1, 2, \ldots$
Die Information wird nicht in einzelnen Niveaus
gespeichert, sondern in der \textbf{Wellenform}
des Zustands --- zum Beispiel als kohärenter Zustand
$\ket{\alpha}$ mit komplexer Amplitude $\alpha$.

Quantenpunkte sind \emph{keine} harmonischen Oszillatoren
(ihre Niveauabstände sind nicht äquidistant: $E_n \sim n^2$) ---
aber \textbf{optische Kavitäten} und
\textbf{Supraleiter-Resonatoren} (Circuit-QED) sind es näherungsweise.

\begin{keypoint}[Diskret oder kontinuierlich --- eine Wahl der Beschreibungsebene]
In FFGFT sind beide Beschreibungen Facetten derselben
Physik. Diskrete Zustände entstehen aus der
endlichen Geometrie des Resonators --- der Resonator
ist endlich, also sind seine Moden abzählbar.
Eine unendliche Modenzahl (idealer harmonischer Oszillator)
ist der Grenzfall eines unendlich tiefen Potentials
mit gleichmäßigen Wänden --- eine Idealisierung.
Reale Systeme sind immer diskret.
Kontinuierliche Variablen sind eine nützliche Näherung
wenn die Modenzahl so groß ist dass man sie
als kontinuierlich behandeln kann ---
analog zur klassischen Physik als Limes vieler Quantenzustände.
\end{keypoint}

% ============================================================
\section{Superposition --- was zwischen den Zuständen liegt}
\label{sec:superposition}
% ============================================================

Bisher: diskrete Zustände $\ket{0}$, $\ket{1}$, \ldots
Das Besondere an Quantensystemen ist nicht dass sie
diskrete Zustände haben --- klassische Münzen auch.
Das Besondere ist die \textbf{Superposition}:

\begin{equation}
  \ket{\psi} \;=\; \alpha\ket{0} + \beta\ket{1},
  \quad |\alpha|^2 + |\beta|^2 = 1.
\end{equation}
Das ist kein Unwissen über den Zustand ---
das ist ein genuiner Zwischenzustand.

Auf der Bloch-Kugel (Abbildung der Qubit-Zustände
auf eine Kugeloberfläche) entspricht jeder Punkt
einem möglichen Quantenzustand ---
unendlich viele kontinuierliche Punkte auf der Kugel,
nicht nur die zwei Pole.

\begin{foundation}[FFGFT: Superposition als intermediäre Torsion]
Spin-up und Spin-down sind die beiden stabilen
Torsionsextreme des Resonators.
Ein Superpositonszustand $\alpha\ket{\uparrow}+\beta\ket{\downarrow}$
ist eine \textbf{metastabile Zwischentorsion} ---
nicht stabil im Grundzustand, aber stabil auf
der Zeitskala der Kohärenzzeit $T_2$.

Die Phase $\phi = \arg(\alpha/\beta)$ ist die
Rotationsstellung dieser Zwischentorsion.
Dekohärenz ist der Prozess bei dem die Umgebung
diese intermediäre Torsion in eine der beiden
stabilen Extremkonfigurationen drängt.
\end{foundation}

Für ein Qudit der Dimension $d$ liegt der
allgemeine Zustand auf einer $(d-1)$-dimensionalen
komplexen Kugel --- dem \textbf{$\mathbb{CP}^{d-1}$}.
Mit wachsendem $d$ wächst der Zustandsraum
exponentiell: $n$ Qudits der Dimension $d$
spannen einen Raum der Dimension $d^n$ auf.

% ============================================================
\section{Dekohärenz --- der Feind der Superposition}
\label{sec:dekohärenz}
% ============================================================

Superposition ist fragil. Die Umgebung misst ständig
mit --- jede Wechselwirkung mit einem Phonon,
einem Kernspin, einem Photon ist eine
(ungewollte) Messung die den Zustand kollabiert.

Zwei Zeitskalen:

\textbf{$T_1$ (Relaxationszeit)}: Zeit bis das System
aus dem angeregten Zustand in den Grundzustand gefallen ist.
Energierelaxation --- der Zustand $\ket{1}$ wird zu $\ket{0}$.

\textbf{$T_2$ (Kohärenzzeit)}: Zeit bis die Phase $\phi$
der Superposition zerstört ist.
Immer $T_2 \leq 2T_1$ (Phasenrelaxation ist mindestens
so schnell wie Energierelaxation, meist schneller).

\begin{center}
\resizebox{\textwidth}{!}{%
\begin{tabular}{lrrl}
\toprule
\textbf{System} & \textbf{$T_1$} & \textbf{$T_2$} & \textbf{Hauptmechanismus} \\
\midrule
Elektron-Spin in GaAs       & $\sim 1$\,ms   & $\sim 10$\,ns   & Kernspin-Bad \\
Elektron-Spin in Si/SiGe    & $> 1$\,s       & $\sim 10$\,ms   & $^{28}$Si (kernspinfrei) \\
Elektron-Spin in P:Si       & $> 1$\,s       & $> 100$\,ms     & Isotopische Reinigung \\
Exziton in CdSe-QD          & $\sim 1$\,ns   & $\sim 10$\,ps   & Phonon-Streuung \\
Supraleiter-Transmon        & $\sim 100$\,µs & $\sim 100$\,µs  & Dielektrische Verluste \\
\bottomrule
\end{tabular}
}
\end{center}

\begin{keyresult}[Si/SiGe: FFGFT-Bedeutung von $T_2 \sim 10$\,ms]
Silizium mit angereichertem $^{28}$Si (kernspinfrei)
eliminiert das Kernspin-Bad fast vollständig.
In FFGFT-Sprache: der Resonator wird von einer
Quelle zufälliger Torsionsstörungen (Kernspins)
befreit. Was übrig bleibt ist die intrinsische
Geometrie des Resonators --- und damit
die fundamentale Grenze $T_2 \leq 2T_1$.
Die lange Kohärenzzeit in $^{28}$Si ist ein
direkter Beleg dafür dass Bedingung~1
(geometrische Perfektion) die entscheidende ist.
\end{keyresult}

% ============================================================
\section{Offene Verbindungen zur FFGFT}
\label{sec:offene-verbindungen}
% ============================================================

Die folgenden Punkte sind noch nicht ausgearbeitet ---
sie markieren Stellen wo das Resonatorbild der FFGFT
konkrete Vorhersagen machen könnte:

\begin{significance}[Offene Fragen für weitere Kapitel]
\textbf{1. Qudit-Niveaustruktur und $\Kfrak$:}\\
Der fraktale Korrekturfaktor $\Kfrak$ (Dokument~173)
verschiebt alle Energieniveaus um denselben relativen Betrag.
Für ein Qudit der Dimension $d$ bedeutet das
$d-1$ verschobene Übergänge --- eine strukturierte
spektrale Signatur die von der idealen Parabolik abweicht.
Ist diese Abweichung messbar mit moderner
Mehrphotonen-Spektroskopie?\\[6pt]

\textbf{2. Spin-Orbital-Kopplung als Torsions-Torsions-Kopplung:}\\
Die Spin-Bahn-Kopplung verbindet zwei geometrische
Freiheitsgrade (Spin = Torsionsrichtung,
Orbital = Torsionsmuster). In FFGFT sollte diese
Kopplung materialabhängig sein --- genau über den
effektiven $g$-Faktor $g^*$ und die
Bandstruktur-Torsion des Kristalls.
Kann $g^*$ aus $\xi$ und der Kristallgeometrie
hergeleitet werden?\\[6pt]

\textbf{3. Dekohärenz-Zeitskalen und $\xi$-Hierarchie --- numerische Analyse:}\\
Gemessene Kohärenzzeiten auf der $\xi$-Stufe $N \approx 8$:
\begin{center}
\begin{tabular}{lll}
\toprule
System & Zeitskala & Verhältnis zu Exziton \\
\midrule
Exziton (CdSe-QD) & $T_2 \sim 10$\,ps & $1$ (Referenz) \\
Spin-Qubit Si/SiGe & $T_2 \sim 10$\,ms & $10^9$ ($= 9$ Größenordnungen) \\
Kernspin P:Si & $T_1 > 1$\,s & $10^{11}$ ($= 11$ Größenordnungen) \\
\bottomrule
\end{tabular}
\end{center}

\textbf{Numerischer Vergleich mit den zwei $\xi$-Parametern:}
\begin{center}
\begin{tabular}{lll}
\toprule
Parameter & $\log_{10}(1/\xi)$ & 2 Stufen $=$ \\
\midrule
$\xipar = 1{,}333 \times 10^{-4}$ (Skalenhierarchie) &
  $3{,}875$ Größenordnungen & $7{,}75$ \\
$\xi_\text{Higgs} = 1{,}038 \times 10^{-5}$ (Kopplungsparameter) &
  $4{,}984$ Größenordnungen & $\mathbf{9{,}97}$ \\
\bottomrule
\end{tabular}
\end{center}

\textbf{Befund:} Mit $\xi_\text{Higgs}$ überbrücken genau \textbf{2 Hierarchiestufen}
exakt $9{,}97$ Größenordnungen --- was zwischen den beiden
gemessenen Verhältnissen ($9$ und $11$ Größenordnungen) liegt
und auf $< 5\,\%$ mit dem Si/SiGe-Wert übereinstimmt.

\textbf{Physikalische Deutung:}
Der Exziton-Resonator auf $\xi$-Stufe $N \approx 8$ koppelt
\emph{direkt} an seine Umgebung (Phononen derselben Stufe):
Dekohärenzrate $\Gamma_\text{Exziton} \propto \xi^0 \cdot \omega$.
Der Spin-Qubit auf derselben Längenskala koppelt an ein
\emph{entferntes} Bad (Kernspins auf Stufe $N \approx 7$):
$\Gamma_\text{Spin} \propto \xi_\text{Higgs}^2 \cdot \omega$ ---
zwei Stufen schwächer, also $T_2 \propto 1/\xi_\text{Higgs}^2
\approx 10^{10}$ mal länger.

\textbf{Schluss und Erklärung:}
Die Übereinstimmung von 2 Hierarchiestufen mit $\xi_\text{Higgs}$ und
dem gemessenen T$_2$-Verhältnis ist kein Zufall ---
sie folgt aus der \textbf{Natur der zwei FFGFT-Räume}:

$\xi_\text{Higgs}$ wirkt im \emph{Zahlenraum / algebraischen Feldraum}
und beschreibt Kopplungen zwischen Quantenzuständen.
Die Dekohärenz ist ein Prozess im Zustandsraum ---
ein Übergang zwischen Quantenzuständen, verursacht durch
Kopplung an die Umgebung.
Deshalb ist $\xi_\text{Higgs}$ der richtige Parameter für Dekohärenzraten.

$\xipar$ wirkt im \emph{geometrischen 3D-Raum}
und beschreibt Längenshierarchien.
Kohärenzzeiten sind keine geometrischen Längen ---
sie sind algebraische Größen im Zustandsraum.
Es wäre physikalisch falsch, $\xipar$ für Dekohärenzraten zu verwenden.

Die zwei $\xi$-Parameter sind \textbf{kein Widerspruch},
sondern komplementäre Beschreibungen verschiedener Aspekte der FFGFT:
Geometrie des Raums ($\xipar$) und Algebra des Zustandsraums ($\xi_\text{Higgs}$).\\[6pt]

\textbf{4. Verschränkung als Torsions-Topologie:}\\
Zwei verschränkte Qubits befinden sich in einem Zustand
der nicht als Produkt zweier Einzelzustände geschrieben
werden kann --- eine gemeinsame Torsionskonfiguration
über räumlich getrennte Resonatoren.
Was bedeutet das für die lokale Topologie
des FFGFT-Feldes?
Gibt es eine natürliche Beschreibung von Bell-Zuständen
als toroidale Feldkonfigurationen?
\end{significance}

% ============================================================
\section{Zusammenfassung: Vom Spin zum Qudit}
\label{sec:zusammenfassung-174}
% ============================================================

\begin{conclusionBox}[Fazit]
Spin-up und Spin-down sind das Minimum ---
zwei Torsionsorientierungen in einem Resonator.

Beeinflussen: Magnetfeld (global), elektrisches Feld
via Spin-Bahn-Kopplung (lokal, schnell), zirkulares Licht
(drehimpulsübertragend), Austauschwechselwirkung (Zwei-QD-Gatter).

Auslesen: Spin-zu-Ladungs-Konversion (praktisch, präzise),
Pauli-Blockade (geometrisch selbst-auslesen),
Fluoreszenz-Polarisation (invasiv, zeitaufgelöst),
\textbf{Faraday/Kerr-Rotation (QND --- zustandserhaltend, wiederholbar)}.

Die QND-Messung ist dabei der physikalisch wichtigste Punkt:
Sie beweist dass der Spinzustand vor der Messung
einen definierten Wert hatte --- die Torsionskonfiguration
ist real und stabil, unabhängig von der Beobachtung.
Das ist der empirische Befund für Determinismus.

Jenseits von zwei Zuständen:
Energieniveau-Qudits, Orbital-Qubits, Exziton-Qudit,
Spin-Orbital-Kombination --- alle in demselben
winzigen Resonator, alle auf der $\xi$-Stufe $N \approx 8$.

Superposition liegt zwischen den Zuständen ---
nicht Unwissen, sondern genuiner Zwischenzustand
(intermediäre Torsion).
Dekohärenz ist ihr natürlicher Feind:
die Umgebung kollabiert die Torsion in eine der
stabilen Extremkonfigurationen.

Geometrische Perfektion (Bedingung~1) ist
die universelle Grundlage --- für jeden dieser Effekte.
Ohne präzisen Resonator kein Qudit, kein Qubit,
kein Spin-Gatter. Die FFGFT-Deutung:
alle diese Phänomene sind verschiedene Ausdrucksformen
stabiler Torsionskonfigurationen in geometrisch
präzisen Hohlräumen auf der Skala $N \approx 8$.
\end{conclusionBox}

% ============================================================
% ============================================================

\begin{thebibliography}{99}

% --- Spin-Manipulation: Magnetfeld / Zeeman ---
\bibitem{Loss1998}
Loss, D. \& DiVincenzo, D.\,P.
\textit{Quantum computation with quantum dots.}
Phys.\ Rev.\ A \textbf{57}, 120--126 (1998).

% --- EDSR ---
\bibitem{Koppens2006}
Koppens, F.\,H.\,L. et al.
\textit{Driven coherent oscillations of a single electron spin in a quantum dot.}
Nature \textbf{442}, 766--771 (2006).

\bibitem{Nowack2007}
Nowack, K.\,C. et al.
\textit{Coherent control of a single electron spin with electric fields.}
Science \textbf{318}, 1430--1433 (2007).

\bibitem{Corna2018}
Corna, A. et al.
\textit{Electrically driven electron spin resonance mediated by
spin--valley--orbit coupling in a silicon quantum dot.}
npj Quantum Inf.\ \textbf{4}, 6 (2018).

% --- Optische Spinpräparation und Manipulation ---
\bibitem{Atature2006}
Atatüre, M. et al.
\textit{Quantum-dot spin-state preparation with near-unity fidelity.}
Science \textbf{312}, 551--553 (2006).

\bibitem{Press2008}
Press, D. et al.
\textit{Complete quantum control of a single quantum dot spin using
ultrafast optical pulses.}
Nature \textbf{456}, 218--221 (2008).

% --- Austauschwechselwirkung ---
\bibitem{Petta2005}
Petta, J.\,R. et al.
\textit{Coherent manipulation of coupled electron spins in semiconductor
quantum dots.}
Science \textbf{309}, 2180--2184 (2005).

% --- Spin-zu-Ladungs-Konversion (Elzerman-Auslese) ---
\bibitem{Elzerman2004}
Elzerman, J.\,M. et al.
\textit{Single-shot read-out of an individual electron spin in a quantum dot.}
Nature \textbf{430}, 431--435 (2004).

% --- Pauli-Blockade ---
\bibitem{Ono2002}
Ono, K. et al.
\textit{Current rectification by Pauli exclusion in a weakly coupled
double quantum dot system.}
Science \textbf{297}, 1313--1317 (2002).

% --- Faraday-/Kerr-Rotation: QND-Auslese ---
\bibitem{Atature2007}
Atatüre, M. et al.
\textit{Observation of Faraday rotation from a single confined spin.}
Nature Phys.\ \textbf{3}, 101--105 (2007).

\bibitem{Berezovsky2008}
Berezovsky, J. et al.
\textit{Picosecond coherent optical manipulation of a single electron spin in a
quantum dot.}
Science \textbf{320}, 349--352 (2008).

\bibitem{Grangier1998}
Grangier, P., Levenson, J.\,A. \& Poizat, J.\,P.
\textit{Quantum non-demolition measurements in optics.}
Nature \textbf{396}, 537--542 (1998).

% --- QND-Messung eines Elektronenspins ---
\bibitem{Delbecq2019}
Delbecq, M.\,R. et al.
\textit{Quantum non-demolition measurement of an electron spin qubit.}
Nature Nanotechnol.\ \textbf{14}, 446--450 (2019).

\bibitem{Xue2020}
Xue, X. et al.
\textit{Quantum non-demolition readout of an electron spin in silicon.}
Nature Commun.\ \textbf{11}, 1178 (2020).

% --- Si/SiGe Kohärenzzeiten ---
\bibitem{Veldhorst2014}
Veldhorst, M. et al.
\textit{An addressable quantum dot qubit with fault-tolerant control-fidelity.}
Nature Nanotechnol.\ \textbf{9}, 981--985 (2014).

\end{thebibliography}

	
\end{document}